\section*{Code and data availability}
All code, result files, pre-registrations and the source of this paper are at \url{https://github.com/arhancanli/null-zoo} (MIT licence); version 1 of the paper is archived at \url{https://doi.org/10.5281/zenodo.23018988}. The JavaScript benchmark in \nolinkurl{src/} is copied unchanged from \url{https://github.com/arhancanli/canlicapital} at commits \nolinkurl{\nz{src.commit.v0}} (v0) and \nolinkurl{\nz{src.commit.v1}} (v1, v1b), where it is maintained. The second implementation, its instrument tests, the v2 and v2b pre-registrations, the run and scoring scripts, the empirical illustration and their outputs are in \nolinkurl{analysis/}. Every result in the text and tables is generated from the result files by \nolinkurl{analysis/tables.mjs}, and the paper reads each one by key; design constants and figures quoted from the literature are typed.

\paragraph{Pre-registration record.} Each pre-registration was committed and pushed before its run (commit, UTC time): v1 \nolinkurl{\nz{prov.v1.prereg_commit}}, \nz{prov.v1.prereg_utc}, results \nolinkurl{\nz{prov.v1.results_commit}}, \nz{prov.v1.results_utc}; v1b \nolinkurl{\nz{prov.v1b.prereg_commit}}, \nz{prov.v1b.prereg_utc}, results \nolinkurl{\nz{prov.v1b.results_commit}}, \nz{prov.v1b.results_utc} (both in \url{https://github.com/arhancanli/canlicapital/pull/379}); v2 \nolinkurl{\nz{prov.v2.prereg_commit}}, \nz{prov.v2.prereg_utc}, run started \nz{prov.v2.run_start_utc}; v2b \nolinkurl{\nz{prov.v2b.prereg_commit}}, \nz{prov.v2b.prereg_utc}, data downloaded \nz{prov.v2b.data_download_utc}, run started \nz{prov.v2b.run_start_utc}. The times are those recorded by git and GitHub; no third-party timestamp was obtained.

\paragraph{Reproduction.} The published v0 result file (SHA-256 \nolinkurl{\nz{v0.sha}}) reproduces byte for byte. Continuous integration checks that reproduction, the instrument tests of both implementations, and that the committed evaluations and tables match the result files; it does not re-run v1, v1b, v2 or v2b, which take hours. We re-ran v1 with \nz{prov.rerun.v1.workers} worker threads instead of \nz{prov.rerun.v1.published_workers}, and v2 with \nz{prov.rerun.v2.workers} worker processes instead of \nz{prov.rerun.v2.published_workers}; both result files were byte-identical (SHA-256 \nolinkurl{\nz{prov.rerun.v1.sha}} and \nolinkurl{\nz{prov.rerun.v2.sha}}). The Python runs used numpy \nz{v2.numpy}; numpy does not promise identical random streams across versions, so a byte-for-byte re-run needs that version. The empirical illustration downloads its data from the French Data Library, whose files are revised over time; the script refuses a file whose SHA-256 differs from the one recorded (\nolinkurl{\nz{emp.zip_sha}}, downloaded \nz{emp.downloaded}).

\section*{Competing interests}
The author designed the luck-equivalent-trials statistic scored here. The author's firm, Canli Capital, publishes a validation API, free at the time of writing, and MCP servers that implement the luck-equivalent-trials test, the haircut, the deflated Sharpe ratio and the data-snooping tests; the v1 and v1b runs set their defaults (the Reality Check for searches, LET-Lo for single series), and the firm may benefit from wider use of these tools. No external funding was received.
